\section{商空间的Hausdorff性}

记号约定:$X$是拓扑空间,$R$是$X$上的等价关系,$\varphi:X\to X/R$是典范满射,相对$R$的二元关系是$C$.

\begin{definition}{Hausdorff等价关系}{}
	若$X/R$是Hausdorff空间,则称$R$是Hausdorff等价关系.
\end{definition}

\begin{remark}{}{}
	\begin{enumerate}
		\item 若$X/R$是Hausdorff空间,则$X$中每一点$x$相对$R$的等价类都是$X$中的闭集.但是,这个条件不是充分的.
		\item $X/R$是Hausdorff空间$\iff$对$X$中任意满足$\left[x\right]_{R}\neq\left[y\right]_{R}$的两点$x,y$,在$X$中有两个关于$R$饱和且不相交的开集,分别包含$\left[x\right]_{R}$和$\left[y\right]_{R}$.
	\end{enumerate}
\end{remark}

\begin{proposition}{相对等价关系的二元关系的闭性和商空间的Hausdorff性}{}
	\begin{enumerate}
		\item 若$X/R$是Hausdorff空间,则$C$是$X\times X$的闭集.
		\item 若$R$是开等价关系,$C$是$X\times X$的闭集,则$X/R$是Hausdorff空间.
	\end{enumerate}
\end{proposition}

\begin{remark}{}{}
	如果$R$不是开等价关系,即便$C$是$X\times X$的闭集,$X/R$也不一定是Hausdorff空间.
\end{remark}

\begin{proposition}{用局部连续映射判断商空间的Hausdorff性}{}
	若对$X$中任意两个不同等价类$M,N$,在$X$中存在包含$M,N$的饱和开集$A$,存在Hausdorff空间$Y$,$f\in\Hom_{\mathsf{Top}}\left(A,Y\right)$,使
	\begin{enumerate}
		\item $f$在$A$中的每个等价类上都取常值,
		\item $f\left(M\right)\neq f\left(N\right)$,
	\end{enumerate}
	则$X/R$是Hausdorff空间.
\end{proposition}

\begin{proposition}{连续映射导出的商空间的Hausdorff性}{}
	设$Y$是Hausdorff空间,$f\in\Hom_{\mathsf{Top}}\left(X,Y\right)$.若$R$是$f$诱导的等价关系,则$X/R$是Hausdorff空间.
\end{proposition}

\begin{proposition}{连续截面与商空间的Hausdorff性}{}
	设$X$是Hausdorff空间.若$\varphi$有连续的右逆,则$X/R$是Hausdorff空间,$\im\left(s\right)$是$X$的闭集.
\end{proposition}






























